% Part: computability
% Chapter: recursive-functions
% Section: non-pr-functions

\documentclass[../../../include/open-logic-section]{subfiles}

\begin{document}

% SEGMENT OLP-0223-S01

\olfileid{cmp}{rec}{npr}
\olsection{முதன்மை மீள்வரையறைக்குரியவை அல்லாத சார்புகள்}

உள்ளுணர்வாகக் கணிக்கத்தக்க அனைத்து சார்புகளையும் முதன்மை
மீள்வரையறைச் சார்புகள் தீர்த்துவிடுவதில்லை. அனைத்து ஓரிட முதன்மை
மீள்வரையறைச் சார்புகளையும் $f_0$, $f_1$, $f_2$,~\dots என,
உள்ளீடு $y$ இல் $f_x$ இன் மதிப்பை விளைவுறக் கணிக்கக்கூடியவாறு,
ஒரு பட்டியலாக்கலாம் என்பது உள்ளுணர்வாகத் தெளிவாக இருக்க வேண்டும்.
வேறுவிதமாகக் கூறினால், பின்வருமாறு வரையறுக்கப்படும் சார்பு $g(x,y)$
கணிக்கத்தக்கது:
\[
g(x,y) = f_x(y)
\]
ஆனால் அப்போது பின்வரும் சார்பும் கணிக்கத்தக்கது:
\begin{eqnarray*}
h(x) & = & g(x,x) + 1 \\
& = & f_x(x) +1.
\end{eqnarray*}
ஒவ்வொரு முதன்மை மீள்வரையறைச் சார்பு $f_i$ க்கும், $h$ மற்றும்
$f_i$ இன் மதிப்புகள் $i$ இல் வேறுபடுகின்றன. எனவே $h$
கணிக்கத்தக்கது; ஆனால் முதன்மை மீள்வரையறைச் சார்பு அல்ல;
$g$ குறித்தும் இதையே கூறலாம். இது கான்டரின் மூலைவிட்டமாக்கல்
வாதத்தின் “விளைவுறு” வடிவம்.

முதன்மை மீள்வரையறைக்குரியவை அல்லாத கணிக்கத்தக்க சார்புகளுக்கு மேலும்
வெளிப்படையான எடுத்துக்காட்டுகளை வழங்கலாம். எடுத்துக்காட்டாக,
$g^n(x)$ என்ற குறிமுறை, மொத்தமாக $n$ எண்ணிக்கையிலான $g$ களைக்
கொண்ட $g(g(\dots g(x)))$ ஐக் குறிக்கட்டும்; மேலும் சார்புகளின்
தொடர் $g_0,g_1,\dots$ ஐப் பின்வருமாறு வரையறுக்க:
\begin{eqnarray*}
g_0(x) & = & x+1 \\
g_{n + 1}(x) & = & g_n^x(x)
\end{eqnarray*}
ஒவ்வொரு சார்பு $g_n$ உம் முதன்மை மீள்வரையறைச் சார்பு என்பதை
உறுதிசெய்யலாம். அடுத்தடுத்து வரும் ஒவ்வொரு சார்பும் அதற்கு
முந்தியதைவிட மிக வேகமாக வளர்கிறது; $g_1(x)$ ஆனது $2x$ க்குச் சமம்;
$g_2(x)$ ஆனது $2^x \cdot x$ க்குச் சமம்; மேலும் $g_3(x)$ ஆனது
$x$ எண்ணிக்கையிலான $2$ களைக் கொண்ட அடுக்குக் கோபுரம் போலத்
தோராயமாக வளர்கிறது. அக்கர்மன்--P\'eter சார்பு சாரத்தில்
$G(x) = g_x(x)$ என்ற சார்பு; இது எந்த முதன்மை மீள்வரையறைச்
சார்பையும்விட வேகமாக வளர்கிறது என்று காட்டலாம்.

முதன்மை மீள்வரையறைச் சார்புகளை எண்ணிடும் சிக்கலுக்குத் திரும்புவோம்.
ஒவ்வொரு முதன்மை மீள்வரையறைச் சார்புக்கும் குறியீட்டுக் குறிமுறைகளை
ஒதுக்கியுள்ளோம் என்பதை நினைவுகூர்க; எனவே குறிமுறைகளை எண்ணிட்டால்
போதும். ஒவ்வொரு குறிமுறை $F$ க்கும் இயல் எண் $\#(F)$ ஒன்றைப்
பின்வருமாறு மீள்வரையறை முறையில் ஒதுக்கலாம்:
\begin{eqnarray*}
\#(0) & = & \langle 0 \rangle \\
\#(S) & = & \langle 1 \rangle \\
\#(\Proj{n}{i}) & = & \langle 2, n, i \rangle \\
\#(\fn{Comp}_{k,l}[H,G_0,\dots,G_{k-1}]) & = & \langle
3,k,l,\#(H),\#(G_0),\dots,\#(G_{k-1}) \rangle \\
\#(\fn{Rec}_l[G,H]) & = & \langle 4, l, \#(G), \#(H) \rangle
\end{eqnarray*}
முந்தைய பிரிவின் குறியீடுகளைப் பயன்படுத்தி, எண்களின் ஒவ்வொரு தொடரையும்
ஓர் இயல் எண்ணாகக் கருதலாம் என்ற உண்மையை இங்கே பயன்படுத்துகிறோம்.
இதன் விளைவாக ஒவ்வொரு குறியீட்டுக்கும் ஓர் இயல் எண் ஒதுக்கப்படுகிறது.
சில தொடர்கள், எனவே சில எண்கள், குறிமுறைகளுக்கு ஒத்திருக்காது என்பது
உண்மை; ஆனால் $i$ அத்தகைய ஒரு குறிமுறையைக் குறியீடாக்கினால்,
குறியீடு $i$ உடைய குறிமுறைக்கான ஓரிட முதன்மை மீள்வரையறைச் சார்பு
$f_i$ ஆகட்டும்; இல்லையெனில் அது மாறிலி $0$ சார்பாகட்டும்.
மொத்த விளைவாக, ஓரிட முதன்மை மீள்வரையறைச் சார்புகளை எண்ணிடுவதற்கான
வெளிப்படையான வழி நமக்குக் கிடைக்கிறது.

(உண்மையில், மாறிலிச் சுழியச் சார்பு போன்ற சில சார்புகள் பட்டியலில்
ஒன்றுக்கு மேற்பட்ட முறை தோன்றும். இது நமது குறியீடாக்கத்தின்
விளைவு மட்டுமல்ல; மாறிலிச் சுழியச் சார்பு ஒன்றுக்கு மேற்பட்ட
குறிமுறைகளைக் கொண்டிருப்பதன் விளைவுமாகும். இம்மீள்தோற்றங்களைக்
கணிக்கத்தக்க வகையில் தவிர்க்க முடியாது என்பதைப் பின்னர் காண்போம்;
எடுத்துக்காட்டாக, கொடுக்கப்பட்ட குறிமுறை மாறிலிச் சுழியச் சார்பைக்
குறிக்கிறதா இல்லையா என்று தீர்மானிக்கும் கணிக்கத்தக்க சார்பு இல்லை.)

இப்போது $g(x,y)$ என்ற சார்பை $f_x(y)$ ஆல் தரலாம்; இங்கே $f_x$
என்பது நாம் இப்போது விவரித்த எண்ணிடலைச் சுட்டுகிறது. $g(x,y)$
கணிக்கத்தக்கது என்பது நமக்கு எவ்வாறு தெரியும்? உள்ளுணர்வாக இது
தெளிவு: $g(x,y)$ ஐக் கணிக்க, முதலில் $x$ ஐ “பிரித்தெடுத்து”,
அது ஓரிடச் சார்புக்கான குறிமுறையா என்று பார்க்க. அவ்வாறெனில்,
உள்ளீடு~$y$ இல் அச்சார்பின் மதிப்பைக் கணிக்க.

\begin{tagblock}{TMs}
\begin{digress}
இதைச் செய்யும் நிரல் ஒன்றை (எடுத்துக்காட்டாக, Java அல்லது C++ இல்)
எழுத முடியும் என்று நீங்கள் ஏற்கெனவே நம்பியிருக்கலாம்---அதற்குச்
சிறிது வேலை தேவைப்படும்! இப்போது உள்ளுணர்வாகக் கணிக்கத்தக்க எதையும்
ஒரு டியூரிங் பொறியால் கணிக்கலாம் என்று கூறும் சர்ச்--டியூரிங்
முன்வைப்பைப் பயன்படுத்தலாம்.

$g(x,y)$ கணிக்கத்தக்கது என்று காட்டுவதற்கான மேலும் நேரடியான வழி,
அதைக் கணிக்கும் டியூரிங் பொறியை வெளிப்படையாக விவரிப்பதாகும்.
இது குறிப்பாக சர்ச்--டியூரிங் முன்வைப்பையும் உள்ளுணர்வுக்கான
முறையீடுகளையும் தவிர்க்கும். டியூரிங் பொறியில்
\emph{உருவகப்படுத்தக்கூடிய} கணக்கீட்டு மாதிரி ஒன்றைப் பயன்படுத்தி
$g(x,y)$ கணிக்கத்தக்கது என்று காட்டுவதற்குப் போதுமான கருவிகளை விரைவில்
உருவாக்கியிருப்போம்; அந்த மாதிரி மீள்வரையறைச் சார்புகள்.
\end{digress}
\end{tagblock}

\end{document}
